\documentclass[10pt,twocolumn]{article}
\usepackage[T1]{fontenc}
\usepackage[utf8]{inputenc}
\usepackage{lmodern}
\usepackage[margin=1.8cm,top=2.4cm,bottom=2cm,columnsep=0.6cm]{geometry}
\usepackage{fancyhdr}
\usepackage{titlesec}
\usepackage{booktabs}
\usepackage{amsmath,amssymb,amsfonts}
\usepackage{microtype}
\usepackage{xcolor}
\usepackage{mdframed}
\usepackage{parskip}
\usepackage{enumitem}
\usepackage{hyperref}

\definecolor{rulered}{RGB}{180,30,30}
\definecolor{boxgray}{RGB}{245,245,245}
\definecolor{keywordgray}{RGB}{100,100,100}

\pagestyle{fancy}
\fancyhf{}
\fancyhead[L]{\textsc{\small Claude Mag}}
\fancyhead[C]{\small\textit{Machine Learning Research Digest}}
\fancyhead[R]{\small 2026-08-10}
\fancyfoot[C]{\small\thepage}
\renewcommand{\headrulewidth}{0.8pt}

\titleformat{\section}{\normalsize\bfseries\color{rulered}}{}{0pt}{}%
  [\vspace{-4pt}\color{rulered}\rule{\columnwidth}{0.4pt}]
\titlespacing{\section}{0pt}{8pt}{4pt}

\hypersetup{colorlinks=true,urlcolor=rulered,linkcolor=black}

\begin{document}
\setlength{\parindent}{0pt}
\setlength{\parskip}{4pt}

{\small\color{keywordgray}\textbf{Keywords:}
  post-training quantization \textbullet{}
  multi-precision LLM \textbullet{}
  residual quantization \textbullet{}
  efficient inference}\par\vspace{4pt}

{\LARGE\bfseries\raggedright
  Recurrent Residual Quantization: A Progressive
  Multi-Precision Representation for LLMs}\par\vspace{4pt}

{\large\itshape\raggedright
  One Checkpoint, Four Bit-Widths: A Calibration-Free
  Framework that Builds 2-, 4-, 6-, and 8-Bit Models
  from a Single Residual Stack}\par\vspace{6pt}

{\small\textbf{By} Yu Luo, Bo Dong, Wenhua Cheng, Haihao Shen}\par\vspace{2pt}

{\small\color{keywordgray}arXiv:2608.04048 \textbullet{} NeurIPS 2026}
\par\vspace{2pt}\noindent{\color{rulered}\rule{\columnwidth}{0.6pt}}\vspace{6pt}

Deploying LLMs across heterogeneous hardware---from cloud GPUs to edge devices---today
requires a separate quantized checkpoint for every target bit-width. Recurrent Residual
Quantization (RRQ) eliminates this overhead by representing weights as a base plus a
sequence of 2-bit residual corrections: reading only the base gives 2-bit inference;
reading base plus $k$ residuals gives $(2+2k)$-bit inference. The full 2/4/6/8-bit
package for Qwen3-8B is constructed in 1{,}293 seconds, $3.3\times$ faster than
MatGPTQ, with no calibration data required.

\begin{mdframed}[backgroundcolor=boxgray,linecolor=rulered,linewidth=1pt,
    innerleftmargin=6pt,innerrightmargin=6pt,
    innertopmargin=6pt,innerbottommargin=6pt]
\textbf{\small AT A GLANCE}\par\vspace{4pt}
\begin{description}[leftmargin=0pt,labelsep=4pt,itemsep=2pt,parsep=0pt,
    font=\small\bfseries]
  \item[Problem:] Each target bit-width requires a separate quantization run and
    checkpoint, multiplying storage and maintenance costs linearly with the number
    of deployment targets.
  \item[Why it matters:] A single multi-precision checkpoint enables hardware-adaptive
    serving: the same model file serves data-center GPUs at 8-bit and edge CPUs at 2-bit.
  \item[Method:] Progressive 2-bit residuals applied to the quantization error of the
    previous step; no calibration data; no joint multi-bit optimization.
  \item[Result:] $3.3\times$ faster construction than MatGPTQ; competitive perplexity
    at all four precisions; single-file deployment.
  \item[Evidence:] Benchmarked on Qwen3-8B and LLaMA-3-8B; accepted to NeurIPS 2026.
\end{description}
\end{mdframed}

\section{The Big Picture}

Large language models are increasingly deployed across a spectrum of hardware: high-memory
data-center GPUs running 8- or 16-bit inference, on-device accelerators requiring 4-bit,
and low-power edge CPUs that may only support 2-bit integer arithmetic. Current practice
requires a separate post-training quantization (PTQ) run---often including an expensive
calibration step---for each target precision. A team deploying Llama-3 at 2-, 4-, and
8-bit must run three separate GPTQ pipelines and store three separate checkpoints.

Recurrent Residual Quantization (RRQ) collapses this into one. A single RRQ construction
run produces a structured checkpoint from which any of the four target precisions can be
read at inference time by selecting how many correction layers to load. There is no
trade-off in quality: each precision performs comparably to a dedicated GPTQ checkpoint
optimized for that bit-width.

\section{Why Existing Approaches Fall Short}

\textbf{Calibration-based PTQ (GPTQ, AWQ, MatGPTQ)} achieves near-FP16 quality at
4-bit and above, but requires a calibration dataset (typically 128 sequences of 2048
tokens) and second-order Hessian computations that dominate construction time. A new
bit-width requires a new run; the runs cannot be composed.

\textbf{Residual Vector Quantization (RVQ)} decomposes weights into a codebook plus
residual codebooks, enabling multiple precisions. However, RVQ optimizes all codebooks
jointly, which is expensive, and the discrete codebook structure is poorly suited to
the per-layer, per-group statistics of LLM weights.

\textbf{Non-uniform quantization (NF4, SpQR)} adapts bit allocation to weight
distributions but is fixed to a single precision per run.

RRQ's key insight is that greedy round-to-nearest residuals---not joint optimization---are
sufficient to close the quality gap between adjacent bit-widths, because each residual
corrects for the error of the previous quantization step independently.

\section{Core Mechanism}

RRQ builds a weight approximation $\hat{w}^{(k)}$ by iteratively quantizing the error
of the previous approximation:
\[
\hat{w}^{(0)} = Q_2(w), \quad
r^{(j)} = w - \hat{w}^{(j-1)}, \quad
\hat{w}^{(j)} = \hat{w}^{(j-1)} + Q_2(r^{(j)}),
\]
where $Q_2(\cdot)$ denotes round-to-nearest 2-bit scalar quantization with per-group
scales computed from the range of the argument. After $k$ iterations, the approximation
has effective precision $(2 + 2k)$ bits. Concretely:
\begin{itemize}[itemsep=2pt,parsep=0pt]
  \item $\hat{w}^{(0)}$: effective 2-bit (base only)
  \item $\hat{w}^{(1)}$: effective 4-bit (base + 1 residual)
  \item $\hat{w}^{(2)}$: effective 6-bit (base + 2 residuals)
  \item $\hat{w}^{(3)}$: effective 8-bit (base + 3 residuals)
\end{itemize}
The checkpoint stores $\hat{w}^{(0)}$ and $r^{(1)}, r^{(2)}, r^{(3)}$ as separate
tensors; inference at precision $(2+2k)$ loads and sums $\hat{w}^{(0)}$ through
$Q_2(r^{(k)})$.

\section{Technical Formulation}

Let $W \in \mathbb{R}^{m \times n}$ be a weight matrix partitioned into groups of
size $g$. For each group $\mathcal{G}_i$, the 2-bit quantizer is:
\[
Q_2(x;\,\mathcal{G}_i) = \mathrm{round}\!\left(
  \frac{x - \min(\mathcal{G}_i)}{\max(\mathcal{G}_i) - \min(\mathcal{G}_i)} \cdot 3
\right) \cdot \frac{\max(\mathcal{G}_i) - \min(\mathcal{G}_i)}{3}
+ \min(\mathcal{G}_i).
\]
The reconstruction error after $k$ residual layers satisfies:
\[
\|w - \hat{w}^{(k)}\|_2^2
= \left\|r^{(k)} - Q_2(r^{(k)})\right\|_2^2
\leq \frac{(\Delta^{(k)})^2}{4},
\]
where $\Delta^{(k)} = \max(r^{(k)}) - \min(r^{(k)})$ is the range of the $k$-th
residual, which shrinks geometrically with $k$ as the residuals concentrate near zero.

\section{Learning or Inference Procedure}

\textbf{Construction.}
For each linear layer, the RRQ construction procedure:
(1) computes the 2-bit base $\hat{w}^{(0)}$ via RTN with per-group scales;
(2) computes residuals $r^{(j)}$ and their 2-bit quantizations sequentially;
(3) stores the base and all residuals. No calibration data is touched; construction
is embarrassingly parallel across layers.

\textbf{Inference.}
A deployment configuration specifies the target precision $b \in \{2, 4, 6, 8\}$.
The runtime loads the base and $(b-2)/2$ residual tensors, sums them, and uses the
result as the weight matrix. Memory footprint is proportional to $b$; switching
precision at inference time requires only loading more or fewer residual tensors,
not re-quantizing.

\textbf{Hardware kernels.}
Because both the base and each residual are 2-bit integer tensors with per-group
scales, existing 2-bit matmul kernels (e.g.\ INT2 CUDA kernels) apply directly.
The multi-residual accumulation adds one extra addition per residual per layer.

\section{What the Guarantee Says}

The greedy residual construction gives a monotone quality improvement: adding each
residual layer can only reduce or maintain the reconstruction error, never increase it.
Formally, $\|w - \hat{w}^{(k)}\|_2 \leq \|w - \hat{w}^{(k-1)}\|_2$ for all $k$.
This ensures that the 4-bit model is always at least as good as the 2-bit model,
and the 6-bit model at least as good as the 4-bit model, from the same checkpoint.

\section{Experimental Findings}

\begin{itemize}[itemsep=2pt,parsep=0pt]
  \item \textbf{Construction time (Qwen3-8B):} 1{,}293 seconds for the full
    2/4/6/8-bit package, $3.3\times$ faster than MatGPTQ (4{,}267 seconds).
  \item \textbf{Perplexity (WikiText-2):} RRQ-4bit matches GPTQ-4bit within
    0.1 perplexity points on Qwen3-8B; RRQ-8bit is within 0.05 of FP16.
  \item \textbf{Zero-shot accuracy (MMLU, HellaSwag):} $<$0.5 point degradation
    at 4-bit vs.\ FP16 baseline on both Qwen3-8B and LLaMA-3-8B.
  \item \textbf{Storage:} The complete 4-precision package is $1.15\times$ the
    size of a single 8-bit GPTQ checkpoint.
\end{itemize}

\section{Ablations and Interpretation}

\textbf{Group size.} Larger groups (128 vs.\ 64) reduce storage but increase 2-bit
error; the residual correction partially compensates, narrowing the quality gap between
group sizes. At 4-bit and above, group size has negligible effect.

\textbf{RTN vs.\ GPTQ base.} Starting from a GPTQ-optimized 2-bit base (rather than
RTN) yields 0.3--0.5 perplexity improvement at 2-bit but does not meaningfully affect
4-bit or higher precision, confirming that the residuals correct for RTN's error.

\textbf{Number of residuals.} Quality saturates after 3 residuals (8-bit effective),
matching the known empirical ceiling of scalar PTQ beyond 8-bit for LLMs.

\section*{Reference}

Yu Luo, Bo Dong, Wenhua Cheng, Haihao Shen.
\textbf{Recurrent Residual Quantization: A Progressive Multi-Precision Representation
for LLMs}. arXiv:2608.04048, NeurIPS 2026.
\url{https://arxiv.org/abs/2608.04048}

\end{document}
